%% Disjoint registry slug for ED26-P07-monotonicity-0001.
\calculussheet{sheet:p07-monotonicity}{Marked monotonicity calculus}{Intrinsic term types have marked arrows; \(\preceq\) orders admissible coercions, \(\approx\) is typed alpha-beta-eta conversion, and assertions \(M\precsim_\tau N\) compare terms whose intrinsic types are below \(\tau\). Every inference is well typed. Premises in \(\Gamma\) are universally interpreted; derivations are finite. Semantic coercions satisfy equation (C) in Definition~\ref{def:p07-monotonicity-prestructure}. Source: \cite[Sections 3--9]{IcardMossTune2017Monotonicity}.}

\hbrule{rule:p07-monotonicity-variable}{Variable}{}{x:\sigma}{\(x\) has declared type \(\sigma\); its sole occurrence has polarity \(+\).}
\hbrule{rule:p07-monotonicity-constant}{Constant}{}{c:\sigma}{\(c\) has declared type \(\sigma\); it has no free-variable occurrences.}
\hbrule{rule:p07-monotonicity-application}{Application}{M:\sigma\xrightarrow m\tau\qquad N:\rho}{M(N):\tau}{\(\rho\preceq\sigma\). Function-part occurrences retain polarity; an argument occurrence of polarity \(n\) receives polarity \(m\circ n\).}
\hbrule{rule:p07-monotonicity-abstraction}{Abstraction}{x:\sigma\qquad M:\tau}{\lambda x.M:\sigma\xrightarrow m\tau}{If there is a free \(x\) and all free \(x\)-occurrences in \(M\) have sign \(+\), choose \(m=+\); if they all have sign \(-\), choose \(m=-\); otherwise, including a vacuous binder, choose \(m=\cdot\). Other free occurrences retain their polarities.}

\hbrule{rule:p07-monotonicity-eq-ref}{Conversion reflexivity}{}{M\approx M}{\(M\) is a well-typed term.}
\hbrule{rule:p07-monotonicity-eq-sym}{Conversion symmetry}{M\approx N}{N\approx M}{\none}
\hbrule{rule:p07-monotonicity-eq-trans}{Conversion transitivity}{M\approx N\qquad N\approx P}{M\approx P}{\none}
\hbrule{rule:p07-monotonicity-alpha}{Alpha}{}{\lambda x.M\approx\lambda y.M[y/x]}{\(y\) has the type of \(x\) and is fresh for \(M\), both free and bound.}
\hbrule{rule:p07-monotonicity-beta}{Beta}{}{(\lambda x.M)N\approx M[N/x]}{The application is well typed; substitution is capture avoiding, with bound variables renamed fresh for \(N\).}
\hbrule{rule:p07-monotonicity-eta}{Eta}{}{\lambda x.M(x)\approx M}{\(x\notin\FV(M)\); the application is well typed, allowing a binder type below the argument type of \(M\). Semantic equality is required at each common comparison type by (E).}
\hbrule{rule:p07-monotonicity-xi}{Xi}{M\approx N}{\lambda x.M\approx\lambda x.N}{All displayed terms are well typed; the binder has the same declared type on both sides.}
\hbrule{rule:p07-monotonicity-congruence}{Cong}{M\approx M'\qquad N\approx N'}{M(N)\approx M'(N')}{Both applications are well typed.}

\hbrule{rule:p07-monotonicity-premise}{Premise}{}{\Gamma\vdash_{\mathrm{mon}}M\precsim_\tau N}{\(M\precsim_\tau N\in\Gamma\).}
\hbrule{rule:p07-monotonicity-ref}{Ref}{}{\Gamma\vdash_{\mathrm{mon}}M\precsim_\tau M}{\(\operatorname{tp}(M)\preceq\tau\).}
\hbrule{rule:p07-monotonicity-trans}{Trans}{\Gamma\vdash_{\mathrm{mon}}M\precsim_\tau N\qquad\Gamma\vdash_{\mathrm{mon}}N\precsim_\tau P}{\Gamma\vdash_{\mathrm{mon}}M\precsim_\tau P}{All three intrinsic types are below \(\tau\).}
\hbrule{rule:p07-monotonicity-point}{Point}{\Gamma\vdash_{\mathrm{mon}}M\precsim_{\sigma\xrightarrow m\tau}N}{\Gamma\vdash_{\mathrm{mon}}M(O)\precsim_\tau N(O)}{\(\operatorname{tp}(O)\preceq\sigma\); both head types are below the displayed comparison arrow, so both applications type-check.}
\hbrule{rule:p07-monotonicity-mono}{Mono}{\Gamma\vdash_{\mathrm{mon}}N\precsim_\sigma O}{\Gamma\vdash_{\mathrm{mon}}M(N)\precsim_\tau M(O)}{\(M\) has intrinsic type \(\sigma\xrightarrow+\tau\).}
\hbrule{rule:p07-monotonicity-anti}{Anti}{\Gamma\vdash_{\mathrm{mon}}N\precsim_\sigma O}{\Gamma\vdash_{\mathrm{mon}}M(O)\precsim_\tau M(N)}{\(M\) has intrinsic type \(\sigma\xrightarrow-\tau\).}
\hbrule{rule:p07-monotonicity-pres}{Pres}{\Gamma\vdash_{\mathrm{mon}}M\precsim_\sigma N}{\Gamma\vdash_{\mathrm{mon}}M\precsim_\tau N}{\(\sigma\preceq\tau\).}
\hbrule{rule:p07-monotonicity-equiv}{Equiv}{}{\Gamma\vdash_{\mathrm{mon}}M\precsim_\tau N}{\(M\approx N\), and both intrinsic types are below \(\tau\).}
\hbrule{rule:p07-monotonicity-func}{Func}{\Gamma\vdash_{\mathrm{mon}}M\precsim_\tau N}{\Gamma\vdash_{\mathrm{mon}}\lambda x.M\precsim_{\sigma\xrightarrow m\tau}\lambda x.N}{\(x:\sigma\); \(m\) is above both abstraction markings. Every member of \(\Gamma\) is universal, so \(x\) need not be fresh for \(\Gamma\).}

\hbrule{rule:p07-monotonicity-erased-par-atom}{Erased parallel atom}{}{a\mathrel{\mathsf{Par}_{\beta\eta}}a}{Auxiliary ordinary unmarked syntax only; \(a\) is a variable or inert constant.}
\hbrule{rule:p07-monotonicity-erased-par-app}{Erased parallel application}{M\mathrel{\mathsf{Par}_{\beta\eta}}M'\qquad N\mathrel{\mathsf{Par}_{\beta\eta}}N'}{MN\mathrel{\mathsf{Par}_{\beta\eta}}M'N'}{Auxiliary ordinary unmarked syntax only.}
\hbrule{rule:p07-monotonicity-erased-par-abs}{Erased parallel abstraction}{M\mathrel{\mathsf{Par}_{\beta\eta}}M'}{\lambda x.M\mathrel{\mathsf{Par}_{\beta\eta}}\lambda x.M'}{Auxiliary ordinary unmarked syntax only; alpha-rename to avoid capture.}
\hbrule{rule:p07-monotonicity-erased-par-beta}{Erased parallel beta}{M\mathrel{\mathsf{Par}_{\beta\eta}}M'\qquad N\mathrel{\mathsf{Par}_{\beta\eta}}N'}{(\lambda x.M)N\mathrel{\mathsf{Par}_{\beta\eta}}M'[N'/x]}{Auxiliary ordinary unmarked syntax only; substitution is capture avoiding.}
\hbrule{rule:p07-monotonicity-erased-par-eta}{Erased parallel eta}{M\mathrel{\mathsf{Par}_{\beta\eta}}M'}{\lambda x.Mx\mathrel{\mathsf{Par}_{\beta\eta}}M'}{Auxiliary ordinary unmarked syntax only; \(x\notin\FV(M)\).}
